\documentclass[conference]{IEEEtran}
\usepackage[T1]{fontenc}
\usepackage[utf8]{inputenc}
\usepackage[brazilian]{babel}
\usepackage{amsmath}
\usepackage{booktabs}
\usepackage[hyphens]{url}
\usepackage{microtype}
\usepackage{cite}

% O resumo fica em ingles, como na versao anterior; babel trocaria para "Resumo".
\addto\captionsbrazilian{\renewcommand{\abstractname}{Abstract}}

\newcommand{\cls}[1]{\textsc{#1}}
\newcommand{\snortimpl}{\ensuremath{\mathrm{Snort}(r,p) \Rightarrow \mathrm{P4\_candidato}(r,p)}}

\begin{document}

\title{Compilação de Cabeçalhos de Regras do Snort~3 em\\
Pré-filtros P4 Conservadores: Avaliação no BMv2}

\author{
\IEEEauthorblockN{Arthur de Freitas Medeiros}
\IEEEauthorblockA{\textit{Universidade Federal do Pampa (UNIPAMPA)}\\
Alegrete, Brazil\\
arthurfreitas.aluno@unipampa.edu.br}
\and
\IEEEauthorblockN{Silvio E. Quincozes}
\IEEEauthorblockA{\textit{Faculty of Computer Science (FACOM)}\\
\textit{Federal University of Uberlândia}\\
Monte Carmelo, Brazil\\
sequincozes@ufu.br}
}

\maketitle

\begin{abstract}
Programmable data planes can evaluate lightweight packet predicates before a
software intrusion detection system (IDS), but Snort~3 rules also depend on
payload, stream, flow, and service semantics that a P4 target cannot
represent. We present a semantics-aware, target-aware Snort-to-P4 compiler
that classifies the relation between each rule and its generated predicate as
\cls{Exact}, \cls{Necessary}, or \cls{Unsupported}, records every omitted
component in a compilation manifest, and fails open for packets outside an
explicitly enforced IPv4 domain. Following the Snort~3.12.2.0 source, it
keeps only the header addresses of 2,931 service-based rules, whose ports and
protocol Snort does not check once it identifies the flow's service; on
BMv2, a plan from an earlier compilation that kept them marked no-match the
witnesses of 190 HTTP rules that alerted off their header port. Among 4,017 evaluated
rules, 3,996 were classified as necessary, 21 as unsupported, and none as
exact; the backend produced 235 P4 entries installed on BMv2. The global
binary prefilter nevertheless marked 100\% of parser-eligible packets as
candidates in three CICIDS2017 traces replayed on BMv2. Observed alerts were
preserved, and an end-to-end campaign over 1,000 sampled rules, with Snort~3
as the oracle and BMv2 as the target, found no violation. Retained
provenance reduced the mean per-packet candidate set to about four of 739
P4-eligible rules, but only by 18.3\% over the full ruleset, against 81.8\%
under the earlier compilation. Semantic safety and operational selectivity
are distinct properties, and what a rule-derived predicate guarantees, and
appears to save, depends on how the IDS evaluates the rule.
\end{abstract}

\section{Introdução}

\textbf{Contexto.} Este trabalho trata da tradução automática de regras de
sistemas de detecção de intrusão (IDS) baseados em assinatura para switches
programáveis. O IDS adotado é o Snort~3~\cite{roesch1999snort,snort3manual},
cujas regras combinam condições de cabeçalho (protocolo, endereços, portas,
flags TCP) com condições que dependem do conteúdo do pacote, do estado do
fluxo e do serviço de aplicação. A linguagem alvo é P4~\cite{bosshart2014p4},
que descreve o processamento de pacotes no plano de dados e executa em
velocidade de linha, mas oferece um modelo computacional restrito: casamento
por tabelas sobre campos de cabeçalho, sem inspeção de payload arbitrária e
sem estado por fluxo no modelo usado. A ideia explorada é a pré-filtragem:
deixar o switch decidir cedo quais pacotes precisam da inspeção completa do
IDS em software.

\textbf{Motivação.} A inspeção em software é o gargalo reconhecido dessa
arquitetura, e a literatura já demonstrou que o plano de dados programável
pode aliviá-la. P4ID~\cite{lewis2019p4id} deriva filtros diretamente de
regras Snort e reporta redução no volume encaminhado ao IDS. O
4MIDable~\cite{lewis2023fourmidable} reporta reduções no tráfego processado
sem alterar as taxas de detecção observadas, em alvos BMv2 e Tofino. O
P4-ONIDS~\cite{tavares2021p4onids} trata explicitamente das restrições de
recursos do plano de dados por meio de agregação e deduplicação de regras. O
conjunto desses resultados estabelece que a abordagem é viável e que há ganho
a obter.

\textbf{Justificativa.} A fraqueza comum a esses trabalhos não está na
viabilidade, e sim no que acontece com a parte da regra que não cabe no alvo.
Uma regra Snort típica conjuga uma condição de cabeçalho com uma condição de
payload ou de estado. Ao gerar o predicado P4, a condição que não cabe precisa
de tratamento explícito, e, no escopo revisado, nenhum dos trabalhos
analisados documenta, regra a regra, qual relação semântica sobrevive à
tradução. Isso importa porque o erro é silencioso e assimétrico: se o
predicado gerado ficar mais \emph{restritivo} que a regra de origem, o
pré-filtro descarta pacotes que a regra original detectaria, e a falha não
aparece em nenhuma métrica de cobertura de tradução. Há ainda uma fonte de
erro menos óbvia, que este trabalho documenta: o cabeçalho de uma regra não
diz, sozinho, o que o IDS exige. No Snort~3, uma regra associada a um serviço
de aplicação deixa de ter porta e protocolo checados quando o serviço do fluxo
é identificado, e um predicado que os exija fica mais restritivo que a regra.

\textbf{Objetivo.} Este trabalho propõe um compilador Snort-para-P4 que torna
explícitas as consequências semânticas da compilação parcial. A propriedade
central que ele busca preservar é
\[
\snortimpl,
\]
ou seja: se a regra $r$ dispararia sobre o pacote $p$, o predicado gerado
obrigatoriamente marca $p$ como candidato. O predicado pode ser mais
permissivo que a regra, nunca mais restritivo. Quando essa garantia não pode
ser assegurada, a regra é rotulada como não suportada em vez de ser traduzida
de forma aproximada, e pacotes que o parser do switch não consegue interpretar
são mantidos como candidatos (\emph{fail-open}). As contribuições específicas
são:
\begin{itemize}
\item Classificar a relação entre cada regra Snort e o predicado P4 gerado em
três categorias --- \cls{Exact}, \cls{Necessary} e \cls{Unsupported} ---
segundo a garantia efetivamente preservada;
\item Registrar em um manifesto de compilação, predicado a predicado, o que
foi preservado, delegado ao Snort, expandido ou tratado por fallback, mantendo
a proveniência entrada-P4 $\rightarrow$ regras de origem através do lowering e
da deduplicação;
\item Implementar o backend completo, da regra até a instalação das entradas
via P4Runtime em um alvo BMv2/v1model, com um domínio de pacotes imposto no
plano de dados e validado por uma suíte controlada e pela contabilidade do
replay;
\item Avaliar separadamente a segurança semântica e a seletividade
operacional, mostrando empiricamente que são propriedades independentes;
\item Executar uma campanha diferencial fim a fim sobre 1.000 regras
amostradas, com o Snort~3 como oráculo e o BMv2 como alvo;
\item Mostrar, a partir do código-fonte do Snort~3.12.2.0, que regras
\emph{service-based} não tornam porta e protocolo do cabeçalho condições
necessárias, compilar as 2.931 regras afetadas segundo essa semântica e medir
no BMv2 o que uma compilação anterior, que os mantinha, deixava escapar;
\item Explorar um uso alternativo da proveniência retida --- em vez de uma
decisão binária, derivar o conjunto de regras candidatas por pacote --- e
medir os limites dessa redução.
\end{itemize}

\section{Fundamentação Teórica}

\textbf{Snort~3 e seu modelo de regras.} O Snort~3 é um IDS de rede baseado em
assinaturas. Cada regra tem um cabeçalho --- ação, protocolo, endereço e porta
de origem, direção, endereço e porta de destino --- e um corpo de opções entre
parênteses. As opções incluem tanto condições de detecção (\texttt{content},
\texttt{pcre}, \texttt{flow}, \texttt{flowbits}, \texttt{flags}) quanto
metadados (\texttt{msg}, \texttt{sid}, \texttt{classtype}). A distinção
relevante aqui é que só uma parte dessas condições é avaliável sobre campos de
cabeçalho; o restante exige inspeção de payload ou estado acumulado de
conexões anteriores.

\textbf{Regras por porta e por serviço.} O Snort~3 distingue regras
\emph{port-based} de regras \emph{service-based}. Uma regra é service-based
quando declara a opção \texttt{service}, usa uma opção específica de um
serviço ou nomeia um serviço no cabeçalho, como em
\texttt{alert http}~\cite{snort3manual}. Numa regra port-based, todo
componente do cabeçalho precisa casar. Numa regra service-based, uma vez que o
Snort identifica o serviço do fluxo como um dos serviços da regra, os
endereços continuam sendo checados, mas as portas e o protocolo não
(Seção~\ref{sec:service}). Um predicado de cabeçalho extraído da regra é,
portanto, uma aproximação cuja relação com a regra precisa ser declarada, e
não presumida.

\textbf{P4 e o modelo v1model.} P4~\cite{p4spec2023} é uma linguagem para
descrever o processamento de pacotes no plano de dados. O programa define o
parser de cabeçalhos, tabelas de casamento e ações. O modelo arquitetural
usado aqui é o v1model. Uma tabela P4 casa campos extraídos do cabeçalho
segundo tipos de casamento (\texttt{exact}, \texttt{lpm}, \texttt{ternary},
\texttt{range}) e dispara uma ação. O modelo não oferece varredura de payload
nem estado por fluxo, o que estabelece o limite do que pode ser traduzido.

\textbf{BMv2 e P4Runtime.} O BMv2~\cite{bmv2} é o switch de referência em
software usado como alvo de execução. As entradas das tabelas não fazem parte
do programa P4: são instaladas em tempo de execução pelo plano de controle via
P4Runtime~\cite{p4runtime2024}, um protocolo gRPC que usa o arquivo P4Info
gerado pelo compilador P4 para identificar tabelas, campos e ações por
identificadores numéricos.

\textbf{Pré-filtragem e predicados necessários.} Um pré-filtro é uma condição
avaliada antes do IDS que decide se o pacote precisa de inspeção. Diz-se que o
predicado $p_r$ é \emph{necessário} para a regra $r$ quando todo pacote que
satisfaz $r$ também satisfaz $p_r$. Um predicado necessário pode ser usado com
segurança para pré-filtragem, porque nunca descarta um pacote que a regra
detectaria --- ainda que retenha tráfego adicional. É essa noção que
fundamenta a propriedade preservada neste trabalho.

\textbf{CICIDS2017.} Conjunto de traços de tráfego produzido pelo Canadian
Institute for Cybersecurity~\cite{sharafaldin2018cicids}, com janelas
rotuladas por tipo de ataque. As janelas de quinta-feira (SQL Injection, XSS e
Brute Force) foram usadas na avaliação.

\section{Trabalhos Relacionados}

\textbf{P4ID~\cite{lewis2019p4id}.} Deriva filtros de pacotes diretamente de
regras Snort, usando endereços IP, portas de transporte, protocolos e flags
TCP, e combina estágios de processamento sem estado e com estado. Sua
principal contribuição é demonstrar que filtragem derivada de regras é viável
no plano de dados programável e que o encaminhamento ao IDS pode ser reduzido.
A limitação central, do ponto de vista deste trabalho, é que a extração de
campos não vem acompanhada, no escopo revisado, da declaração da relação
semântica entre a regra de origem e o filtro gerado: não fica registrado o que
foi descartado ao converter uma regra que também dependia de payload.

\textbf{P4-ONIDS~\cite{tavares2021p4onids}.} Trata explicitamente das
limitações de recursos de planos de dados restritos, aplicando representação
de faixas, agregação de regras, deduplicação e uso de informação de fluxo para
limitar os pacotes encaminhados ao NIDS. A contribuição é mostrar que a
viabilidade em hardware real depende de compressão da representação. A
limitação é que a agregação e a deduplicação são tratadas como otimizações de
ocupação de tabela, sem que o efeito dessas transformações sobre a garantia
semântica de cada regra agregada seja explicitado.

\textbf{4MIDable~\cite{lewis2023fourmidable}.} Propõe um arcabouço para
integrar middleboxes de segurança existentes a uma infraestrutura P4/SDN. Seu
agente P4ID suporta processamento com estado e realimentação para um IDS não
modificado, e foi avaliado tanto em BMv2 quanto em Tofino, reportando redução
de tráfego processado sem afetar as taxas de detecção observadas. A
contribuição é a integração prática com IDS legado e a validação em hardware.
A limitação é que a preservação da detecção é estabelecida empiricamente sobre
os traços avaliados, e não por uma relação por regra que valha
independentemente do tráfego observado.

\textbf{P4NIPS~\cite{senoski2025p4nips}.} Apresenta uma implementação recente
de prevenção de intrusão em P4, estendendo a linha de detecção para ação de
bloqueio no plano de dados. Contribui com uma realização concreta e atual do
conceito. A limitação, para os fins deste trabalho, é a mesma das anteriores:
o foco está na construção do sistema, não na caracterização do que a tradução
preserva ou perde.

\textbf{eRBF --- Brum et al.~\cite{brum2025rethinking}.} Revisita a
pré-filtragem baseada em regras considerando que o IDS pode precisar de
tráfego que não satisfaz diretamente as condições extraídas da regra. A
abordagem incorpora contexto de fluxo, incluindo estabelecimento e
encerramento de conexões TCP e dependências de protocolo selecionadas, e é
avaliada com Snort~3 sobre conjuntos grandes de regras. A contribuição é
mostrar que pré-filtragem agressiva precisa levar em conta o contexto exigido
pelo IDS --- um argumento diretamente alinhado ao deste trabalho. A limitação
é que a avaliação é feita em software sobre traços, sem tratar do mapeamento
dessas semânticas para as restrições de processamento e de recursos de um
alvo P4.

\textbf{RNA~\cite{hagen2025rna}.} Segue direção complementar, orientada a
eventos: analisa eventos selecionados de scripts Zeek e os mapeia para
templates de protocolo e de evento para dispositivos de encaminhamento
programáveis, junto a mecanismos no host. A contribuição é demonstrar a
geração coordenada de funcionalidade distribuída entre host e plano de dados.
A limitação, quanto ao escopo aqui tratado, é que parte de eventos Zeek
selecionados e de templates predefinidos, e não de uma classificação semântica
geral de regras de IDS por assinatura segundo a garantia preservada.

\textbf{Contexto mais amplo.} Surveys sobre fundamentos e aplicações de
P4~\cite{hauser2023p4survey}, sobre aplicações de segurança em switches
P4~\cite{alsabeh2022p4securitysurvey}, sobre IDS em plano de dados
programável~\cite{fida2026pdpids} e trabalhos orientados à
implementação~\cite{mazloum2025p4security} documentam a amplitude do campo e
reforçam que o mapeamento de funcionalidade de segurança para alvos
programáveis é limitado por recursos e por expressividade.

\subsection{Posicionamento}

A Tabela~\ref{tab:rel} compara os sistemas quanto à origem das regras, ao alvo
avaliado, à existência de uma relação semântica declarada e à proveniência por
regra. O padrão que emerge é consistente: os trabalhos demonstram viabilidade,
integração e eficiência de representação, mas tratam a tradução parcial como
um detalhe de implementação. Em nenhum deles, no escopo revisado, é explícita,
por regra, a garantia que o predicado gerado preserva, nem a rastreabilidade
da entrada instalada de volta às regras que a originaram. É essa lacuna que
este trabalho endereça.

\begin{table}[t]
\centering
\caption{Comparação com os trabalhos relacionados. N.E.: não explícito no
escopo revisado, o que não equivale a ausência do recurso.}
\label{tab:rel}
\footnotesize
\setlength{\tabcolsep}{2pt}
\begin{tabular}{lllcc}
\toprule
\textbf{Trabalho} & \textbf{Origem} & \textbf{Alvo} &
\textbf{\shortstack{Semântica\\declarada}} &
\textbf{\shortstack{Proveniência\\por regra}} \\
\midrule
P4ID & Snort & Plano P4 & N.E. & N.E. \\
P4-ONIDS & NIDS & Plano P4 & N.E. & N.E. \\
4MIDable & Agente P4ID & BMv2, Tofino & N.E. & N.E. \\
P4NIPS & N.E. & P4 & N.E. & N.E. \\
eRBF & Snort~3 & Software & N.E. & N.E. \\
RNA & Zeek & P4 + host & N.E. & N.E. \\
\textbf{Este trabalho} & \textbf{Snort~3} & \textbf{BMv2} &
\textbf{Sim (3 classes)} & \textbf{Sim} \\
\bottomrule
\end{tabular}
\end{table}

\section{Proposta}

A proposta é um compilador que traduz regras Snort~3 em entradas de tabela P4
buscando preservar a propriedade \snortimpl{} e tornando explícito tudo aquilo
que não pôde ser traduzido. O pipeline tem nove estágios:
\begin{enumerate}
\item \textbf{Parser} --- lê a sintaxe da regra Snort~3 e produz uma
representação estruturada, rejeitando as formas que o próprio Snort
rejeitaria;
\item \textbf{Resolução de variáveis} --- substitui variáveis de rede
(\texttt{\$HOME\_NET}, \texttt{\$EXTERNAL\_NET}) e de serviço pelos valores do
perfil de configuração;
\item \textbf{IR semântica} --- converte a regra em uma lista de predicados
tipados (cabeçalho, payload, estado de fluxo, metadados);
\item \textbf{Análise de capacidade do alvo} --- para cada predicado, decide se
ele é representável na tabela P4 e atribui um tratamento: \cls{Preserved},
\cls{Expanded}, \cls{Delegated}, \cls{Fallback} ou \cls{Out\_of\_scope};
\item \textbf{Classificador semântico} --- a partir dos tratamentos, classifica
a regra como \cls{Exact} (todo o significado cabe no P4), \cls{Necessary} (o
predicado é mais permissivo, e seguro) ou \cls{Unsupported} (não há predicado
seguro derivável);
\item \textbf{Manifesto} --- registra a decisão de cada predicado de cada
regra, com o motivo, e mantém a ligação entrada $\rightarrow$ regras de
origem;
\item \textbf{Backend P4} --- gera as entradas, expande listas de portas e
regras bidirecionais, normaliza e deduplica;
\item \textbf{Lowering} --- converte para valores e máscaras ternárias;
\item \textbf{Instalação} --- gera o plano P4Runtime e instala no BMv2,
verificando o contrato do P4Info e conferindo a codificação por round-trip.
\end{enumerate}

O ponto de projeto que diferencia a proposta é o tratamento do caso em que a
tradução não é possível. Em vez de descartar a condição que não cabe --- o que
tornaria o predicado mais restritivo e abriria um buraco de detecção ---, a
condição é delegada ao Snort e registrada no manifesto.

\textbf{Domínio de pacotes imposto.} A garantia pressupõe que o plano de dados
leia os campos a que o predicado se refere: um parser que assumisse cabeçalho
IPv4 fixo de 20 bytes leria bytes de opções como portas. O programa P4 define
por isso um domínio suportado $D$: Ethernet sem VLAN com IPv4 direto,
$\mathrm{IHL}=5$, sem fragmentação (flags nulas ou DF) e sem erro de parser.
O parser só extrai a camada de transporte dentro de $D$, e o controle de
ingresso repete a checagem antes de aplicar a tabela. Cada pacote incrementa
exatamente um de três contadores: \texttt{P4\_CANDIDATE} (alguma entrada casou
dentro de $D$), \texttt{P4\_NO\_MATCH} (nenhuma casou dentro de $D$) e
\texttt{PACKET\_DOMAIN\_FALLBACK} (fora de $D$, marcado como candidato
incondicionalmente). O veredito \emph{no-match}, o único que poderia justificar
retirar um pacote do IDS, só é emitido dentro de $D$.

\textbf{Cabeçalho de regras service-based.} Para uma regra service-based, o
compilador mantém só os endereços do cabeçalho e delega porta e protocolo ao
Snort, mantendo o protocolo apenas quando a opção \texttt{flags} já restringe a
regra a TCP, porque essa opção não casa outro protocolo
(Seção~\ref{sec:service}). As regras que nomeiam o serviço no cabeçalho, como
\texttt{alert http}, mantêm as portas do cabeçalho. Uma versão anterior do
compilador tratava toda regra pelo cabeçalho inteiro; seus artefatos estão
arquivados e se regeneram byte a byte.

\section{Avaliação}

A avaliação responde a quatro questões: (Q1) qual a cobertura da compilação
semântica; (Q2) se o backend gerado é viável no alvo; (Q3) qual a seletividade
do pré-filtro resultante; e (Q4) se os alertas observados são preservados e
se a proveniência retida tem utilidade quando a seletividade binária é baixa.
As três primeiras caracterizam a proposta; a última verifica empiricamente a
garantia semântica. A Seção~\ref{sec:service} compara os resultados com os da
compilação anterior.

\subsection{Metodologia}

\textbf{Conjunto de regras e perfil.} O conjunto avaliado é o Snort~3
Community Ruleset, com 4.017 regras, sob o perfil experimental
\texttt{cicids2017}, que resolve \texttt{\$HOME\_NET} para 192.168.10.0/24 e
\texttt{\$HTTP\_PORTS} para 80 e 8080. As classificações dependem dessa
configuração, e o perfil resolvido é registrado junto aos resultados.

\textbf{Tráfego.} Janelas de quinta-feira do CICIDS2017: SQL Injection (35.274
quadros) e XSS (191.848) como traços de desenvolvimento, e Brute Force
(1.108.658) mantido como traço \emph{held-out} --- nenhum ajuste de política
foi feito após inspecioná-lo.

\textbf{Ambiente.} Os alertas de referência vêm do Snort~3.12.2.0. O programa
P4 é compilado pelo \texttt{p4c} 1.2.4.2, com o backend \texttt{p4c-bm2-ss},
e executado no BMv2 (\texttt{simple\_\allowbreak switch\_\allowbreak grpc}
1.15.0, modelo v1model) em Ubuntu 22.04, com as entradas instaladas via
P4Runtime. O replay usa um socket AF\_PACKET que transmite os bytes Ethernet
capturados sem reserialização, a cada 2~ms (20~ms só com os pacotes de
alerta) e com filas veth de 10.000 quadros,
após instalação nova do pipeline, contadores conferidos em zero e uma
auditoria de ociosidade com delta nulo; os quadros enviados são reconciliados
com a soma dos três contadores.

\textbf{Validação do domínio.} Uma suíte controlada envia ao BMv2, com o plano
global de 235 entradas instalado, um quadro por caso e compara o delta dos
contadores com a expectativa offline. Os onze casos cobrem candidatos TCP, UDP
e ICMP, um no-match garantido, flags TCP exatas e mascaradas, opções IPv4,
primeiro fragmento com MF, fragmento não inicial, cabeçalho IPv4 truncado e
cabeçalho TCP truncado.

\textbf{Campanha diferencial fim a fim.} Os traços só exercitam as poucas
regras que disparam neles. Para testar a garantia sobre mais regras, um
gerador sintetiza, a partir da representação intermediária de cada regra, um
fluxo curto cuja testemunha deveria fazer o Snort disparar aquela regra:
handshake TCP quando \texttt{flow} exige sessão estabelecida, cabeçalho com os
valores da regra e cada \texttt{content} no buffer inspecionado. Cada
testemunha tem um vizinho sem o conteúdo (negativo de payload) e, quando o
plano da geração dos casos deixava espaço, um com porta de destino fora de
todas as suas entradas (negativo de cabeçalho). Foi possível sintetizar testemunha para 2.142 das
4.017 regras; das 1.875 excluídas, 1.654 têm opções como \texttt{pcre},
\texttt{byte\_test} ou \texttt{flowbits}, que não se satisfazem escolhendo
bytes. Uma amostra de 1.000 regras, estratificada por protocolo, serviço e
direção com semente fixa, gerou 6.684 quadros, todos dentro de $D$. O Snort lê
o traço com o ruleset completo, o BMv2 recebe os quadros um a um, e cada caso
recebe um veredito: \cls{Confirmado} (o Snort alertou e todo quadro alertado
foi candidato); \cls{Violação} (o Snort alertou e a testemunha, ou outro quadro
que satisfaz o cabeçalho da regra, foi no-match --- o único resultado que
contradiz a propriedade); \cls{Vácuo} (o Snort não alertou, e o caso não
testou nada); ou \cls{Dependência de estado} (a testemunha foi candidata, mas
o Snort atribuiu o alerta a um quadro da sessão que não satisfaz o cabeçalho e
foi no-match). Regras em fallback não são pré-filtradas e são contadas à
parte. O mesmo traço rodou contra o plano global e contra o plano do
discriminador forte (Seção~\ref{sec:resultados}). Como cada testemunha usa
uma porta do próprio cabeçalho, uma segunda campanha leva o fluxo HTTP de 200
regras \texttt{service:http} à porta 18080, que nenhum desses cabeçalhos lista
(Seção~\ref{sec:service}).

\textbf{Métricas.} (i) cobertura de classificação semântica; (ii) número de
entradas instaladas, sua validação e a redução do estado de instalação,
$1 - N_{\mathrm{install}}/N_{\mathrm{source}}$; (iii) taxa de pacotes
candidatos, calculada apenas sobre $D$; (iv) preservação de alertas,
comparando a execução do Snort com o ruleset completo contra a união das
execuções com a partição candidata e com a partição \cls{Unsupported}; e
(v) conjunto de regras candidatas por pacote,
$C(x) = \bigcup_{e \in E(x)} S(e)$, em que $E(x)$ são as entradas que casam o
pacote $x$ e $S(e)$ as regras de origem da entrada $e$. Sua redução é
$R(x) = 1 - |C(x)|/B$ dentro da partição P4-elegível de $B$ regras e
$R_{\mathrm{full}}(x) = 1 - (|C(x)|+F)/4017$ sobre o ruleset, em que $F$ conta
as regras de fallback, que valem para todo pacote.

\textbf{Reprodutibilidade.} O código, o perfil, os scripts de análise e os
artefatos congelados estão em \url{https://github.com/medeiros08/Snort-p4}.
Os dez artefatos do pipeline (manifesto semântico, partição de regras,
entradas P4 e plano P4Runtime) são determinísticos: regenerados em uma cópia
isolada a partir do código versionado, saem idênticos byte a byte, assim como
os da compilação anterior. A suíte automatizada tem 27 módulos de teste; 24
passam e os outros 3, que exigem os bindings P4Runtime/gRPC, só executam no
ambiente com o BMv2.

\subsection{Resultados}
\label{sec:resultados}

\textbf{Cobertura semântica (Q1).} Das 4.017 regras, 3.996 foram classificadas
como \cls{Necessary}, 21 como \cls{Unsupported} e \textbf{nenhuma} como
\cls{Exact}. O parser leu 4.017 de 4.017 regras sem falha. No nível de
predicado: 14.866 preservados, 29.411 delegados ao Snort, 323 expandidos, 65
em fallback e \textbf{zero} não avaliados --- isto é, não houve predicado sobre
o qual o compilador deixasse de decidir. O fato de \cls{Exact} ser zero é o
resultado esperado, e não um defeito: nenhuma regra do conjunto tem todo o seu
significado representável no plano de dados, porque praticamente todas
combinam cabeçalho com \texttt{content} ou com estado de fluxo. As 21
\cls{Unsupported} não são ignoradas: quatro têm cabeçalho não seletivo, dez
são regras de serviço sem cabeçalho tradicional, sete são regras
service-based sem restrição de endereço, cujas únicas condições seletivas
eram porta e protocolo, e todas permanecem no Snort. Das 3.996
\cls{Necessary}, 1.072 têm o predicado derivado do cabeçalho com as portas e 2.924,
service-based, dos endereços (e do TCP, quando há \texttt{flags}).

\textbf{Viabilidade do backend (Q2).} As 4.040 instâncias de predicado
geradas, expandidas e normalizadas resultaram em 231 casamentos únicos e, após
o lowering de flags TCP, em 235 entradas de instalação --- redução de 94,18\%
no \emph{estado da tabela}, não no tráfego. A maior parte vem das regras service-based, cujos
predicados só de endereço coincidem: uma única entrada, destino em
\texttt{HOME\_NET}, carrega a proveniência de 2.043 regras. As 235 entradas
foram codificadas, instaladas no BMv2 e lidas de volta sem erro, dentro da
capacidade de 8.192 entradas da tabela.

\textbf{Domínio de pacotes.} A suíte controlada passou nos onze casos com
deltas exatos; os cinco fragmentados ou truncados caíram em
\texttt{PACKET\_DOMAIN\_FALLBACK}. Nos traços, ficaram fora de $D$ 350, 2.772
e 7.199 quadros, quase todos ARP e IPv6, e todos mantidos como candidatos. O
replay de cada traço inteiro, nos dois planos, contou todos os quadros
enviados, com cada categoria igual à do classificador offline.

\textbf{Seletividade: o resultado negativo (Q3).} O pré-filtro binário global
marcou \textbf{100\% dos pacotes elegíveis} como candidatos nos três traços
(Tabela~\ref{tab:sel}). A análise causal identificou entradas amplas:
\texttt{p4-88} casa qualquer TCP e subsume 162 das 235 entradas, e
\texttt{p4-10} e \texttt{p4-51} casam destino e origem em \texttt{HOME\_NET} e
carregam a proveniência de 2.043 e 900 regras, quase todas service-based. A
ablação com discriminador forte considera forte a entrada que restringe
porta, flags TCP ou prefixo IPv4 diferente de exatamente \texttt{HOME\_NET}, e
mantém uma regra P4-elegível só se todas as suas entradas forem fortes. Ela
removeu dez entradas, manteve 225 para 739 regras elegíveis e pôs 3.257 em
fallback operacional, além das 21 \cls{Unsupported}; mesmo assim, no BMv2,
marcou 92,173--92,298\% dos pacotes nos traços de desenvolvimento e 94,885\%
no held-out. Quatro entradas retidas chegam a 90\% dos pacotes de SQLi e XSS:
uma casa TCP para \texttt{HOME\_NET} com o bit ACK, vinda de uma única regra
SMTP (SID 309, \texttt{flags:A+}), e três casam tráfego de \texttt{HOME\_NET}
para as portas 443, 80 e 53. São discriminadores fortes pelo critério
estrutural, mas coincidem com o tráfego dominante.

\begin{table}[t]
\centering
\caption{Seletividade e preservação por traço, medidas no BMv2 com
contabilidade exata. Candidatos sobre os pacotes em $D$. Alertas do Snort com
o ruleset completo; P4 cand.: pacotes de alerta marcados como candidatos em
cada plano. No plano forte, todo alerta vem de regra em fallback.}
\label{tab:sel}
\footnotesize
\setlength{\tabcolsep}{2pt}
\begin{tabular}{lrrrrrrr}
\toprule
Traço & Em $D$ & \shortstack{Global\\cand.} & \shortstack{Forte\\cand.} &
Alertas & SIDs & \shortstack{P4 cand.\\global} & \shortstack{P4 cand.\\forte} \\
\midrule
SQLi & 34.924 & 100\% & 92,298\% & 150 & 1 & 150 & 0 \\
XSS & 189.076 & 100\% & 92,173\% & 528 & 2 & 528 & 0 \\
Brute Force & 1.101.459 & 100\% & 94,885\% & 2.391 & 5 & 2.391 & 3 \\
\bottomrule
\end{tabular}
\end{table}

\textbf{Preservação de alertas (Q4).} Nos três traços, a partição semântica
não produziu alertas faltantes nem adicionais; as sete regras que a
compilação atual acrescenta às \cls{Unsupported} não alertaram e não
participam de dependência de flowbits, de modo que a medição feita com a
partição anterior vale sem mudança. No BMv2, reproduzindo só os pacotes de
alerta, todos foram candidatos no plano global (Tabela~\ref{tab:sel}). No
plano forte, todos menos três foram no-match, sem perda: todo alerta
observado vem de regra em fallback operacional, que o Snort inspeciona sem o
pré-filtro. A SID 254, responsável por 2.960 dos 3.069 alertas, declara
\texttt{service:dns} e, sem UDP e porta 53, fica só com o destino
\texttt{HOME\_NET}, que não é discriminador forte. Nenhum alerta observado
pertence a regra P4-elegível, e nenhum testa os conjuntos candidatos da
Tabela~\ref{tab:cand}. Registra-se ainda
que nenhuma das 21 regras \cls{Unsupported} disparou e que os traços
acionaram apenas 1, 2 e 5 SIDs distintos, o que limita o poder de
discriminação dessa verificação. Os nomes das janelas também não implicam que
os alertas observados detectaram os ataques rotulados.

\textbf{Campanha fim a fim (Q4).} Nenhum caso resultou em \cls{Violação} em
nenhum dos dois planos (Tabela~\ref{tab:camp}), nenhum quadro caiu fora de
$D$, e as duas atribuições de alertas a casos, por faixa de pacotes e por
timestamp, concordaram em todos. No plano global, 958 regras foram
confirmadas, 40 foram vácuas e duas, Kerberos (SIDs 56583 e 56584), ficaram
em fallback por serem \cls{Unsupported}; as 1.000 testemunhas e os 1.000
negativos de payload foram candidatos, e os 101 negativos de cabeçalho foram
no-match. Nenhuma confirmação desse plano, porém, pode ser atribuída só ao
predicado da própria regra: toda testemunha confirmada também casou uma
entrada de outra regra. No plano do discriminador forte, 134 regras foram confirmadas, 36 delas somente
pelas próprias entradas, de modo que um defeito de lowering nessas regras
teria aparecido como violação; das 826 em fallback operacional, 78 tiveram a
testemunha marcada no-match, o que não viola nada numa regra que o plano não
pré-filtra. Nas duas rodadas, a categoria prevista pelo casador offline
coincidiu com o contador do BMv2 em todos os 2.101 quadros de testemunhas e
vizinhos. Em 606 regras HTTP em fallback, o Snort atribuiu o alerta ao ACK
que dispara a inspeção da mensagem armazenada, e o plano forte marcou 445
desses ACKs como no-match: um sistema que retirasse do IDS os pacotes
no-match dependeria de manter essas regras fora do pré-filtro, como faz o
fallback.

\begin{table}[t]
\centering
\caption{Campanha diferencial sobre a mesma amostra de 1.000 regras, no mesmo
traço. Fallb.: regras que o plano não pré-filtra. ``Só própria'':
confirmadas cuja testemunha casou apenas entradas da própria regra.}
\label{tab:camp}
\footnotesize
\setlength{\tabcolsep}{2.5pt}
\begin{tabular}{lrrrrrrr}
\toprule
Plano & Entr. & Conf. & Viol. & \shortstack{Dep.\\estado} &
Fallb. & Vácuo & \shortstack{Só\\própria} \\
\midrule
Global & 235 & 958 & 0 & 0 & 2 & 40 & 0 \\
Forte & 225 & 134 & 0 & 0 & 826 & 40 & 36 \\
\bottomrule
\end{tabular}
\end{table}

\textbf{Redução do espaço de regras candidatas (Q4).} Usando a proveniência
retida em vez da decisão binária, o conjunto médio de regras candidatas por
pacote, dentro da partição P4-elegível de 739 regras, ficou em 3,97 (SQLi),
3,94 (XSS) e 4,48 (Brute Force \emph{held-out}), redução de 99,39--99,47\%
(Tabela~\ref{tab:cand}). As 3.278 regras de fallback (3.257 operacionais e
21 \cls{Unsupported}) valem para todo pacote; contadas contra o ruleset
completo, as reduções são de 18,30\%, 18,30\% e 18,29\%, logo abaixo do teto
de 18,40\%. Agrupar regras também custa: por protocolo, com um grupo TCP de
642 regras, a redução dentro da partição cai para 20,10--25,38\%; protocolo e
porta da entrada chegam a 98,80--98,96\%.

\begin{table}[t]
\centering
\caption{Conjunto de regras candidatas por pacote (offline, plano do
discriminador forte). $R$ sobre as 739 regras P4-elegíveis;
$R_{\mathrm{full}}$ sobre as 4.017, contando as 3.278 de fallback;
Prot.: redução quando se encaminham grupos inteiros por protocolo.}
\label{tab:cand}
\footnotesize
\setlength{\tabcolsep}{2.5pt}
\begin{tabular}{lrrrrrr}
\toprule
Traço & Média & Mediana & Máx. & $R$ & $R_{\mathrm{full}}$ & Prot. \\
\midrule
SQLi & 3,97 & 3 & 43 & 99,46\% & 18,30\% & 25,38\% \\
XSS & 3,94 & 3 & 43 & 99,47\% & 18,30\% & 23,95\% \\
Brute Force & 4,48 & 3 & 293 & 99,39\% & 18,29\% & 20,10\% \\
\bottomrule
\end{tabular}
\end{table}

\subsection{Regras service-based e a compilação anterior}
\label{sec:service}

A premissa de que uma regra só dispara quando todo o seu cabeçalho casa vale
para o Snort~3 apenas em parte. No Snort~3.12.2.0, antes de checar o
cabeçalho de uma regra com serviços, o Snort compara o serviço identificado
no fluxo com os da regra (\path{src/detection/detection_options.cc} e
\path{src/detection/rtn_checks.cc}). Se o serviço é um deles, as checagens de
porta são puladas, a de protocolo passa incondicionalmente, e só os endereços
são checados; se é outro, a regra é rejeitada; só sem serviço identificado o
cabeçalho inteiro é avaliado. Como a identificação do serviço pode vir da
inspeção do conteúdo, uma requisição HTTP numa porta fora de
\texttt{\$HTTP\_PORTS} ainda é avaliada pelas regras HTTP. A exceção são as
regras que nomeiam o serviço no cabeçalho, como \texttt{alert http}, que o
Snort marca como regras de usuário e que mantêm as portas checadas
(\path{src/parser/parse_rule.cc}).

A versão anterior do compilador tratava o serviço como mais uma condição
delegada e mantinha porta e protocolo no predicado, o que seria correto se o
Snort exigisse serviço \emph{e} cabeçalho. Isso afetava 2.931 regras: 2.928
declaram \texttt{service} (\texttt{http} em 2.314, \texttt{imap} em 150,
\texttt{pop3} em 138, \texttt{smtp} em 131), e três (SIDs 529, 3238 e 3239)
recebem \texttt{dcerpc} e \texttt{netbios-ssn} das opções
\texttt{dce\_iface} e \texttt{dce\_stub\_data}
(\path{src/detection/fp_utils.cc}). Todas tinham protocolo no predicado, e
2.877 também porta. O plano global anterior, de 331 entradas, ainda cobria
quase todas por entradas amplas sem porta, mas o plano do discriminador forte
as removia. O teste com 200 regras \texttt{service:http} na porta 18080
exercita exatamente esse caso: o Snort alertou em 190, e o plano forte
anterior marcou no-match as 190 testemunhas e o ACK ao qual o Snort atribui
cada alerta, isto é, 190 violações. Com a compilação atual, o plano global
marcou as 190 testemunhas como candidatas, e o plano forte tem as 190 regras
em fallback operacional (Tabela~\ref{tab:comp}).

A correção custa sobretudo seletividade: só sete regras perdem cobertura.
Sem porta e protocolo, a
maioria das regras service-based fica só com um endereço em
\texttt{HOME\_NET}, que o discriminador forte não considera discriminante, e
a partição elegível cai de 3.611 para 739 regras. A redução do conjunto
candidato sobre o ruleset, que na compilação anterior chegava a 81,81\% no
held-out, vinha em grande parte de portas que o Snort~3 não checa depois de
identificar o serviço.

\begin{table}[t]
\centering
\caption{Compilação anterior (cabeçalho inteiro) e atual (só endereços nas
regras service-based). Valores de traço no Brute Force held-out.}
\label{tab:comp}
\footnotesize
\setlength{\tabcolsep}{3pt}
\begin{tabular}{lrr}
\toprule
 & Anterior & Atual \\
\midrule
\cls{Necessary} / \cls{Unsupported} & 4.003 / 14 & 3.996 / 21 \\
Entradas (global / forte) & 331 / 322 & 235 / 225 \\
Regras elegíveis no forte & 3.611 & 739 \\
Candidatos no forte & 98,942\% & 94,885\% \\
$R_{\mathrm{full}}$ & 81,81\% & 18,29\% \\
Fora da porta: violações / fallback & 190 / 0 & 0 / 190 \\
\bottomrule
\end{tabular}
\end{table}

\subsection{Discussão dos resultados}

Os achados centrais são três. O primeiro é negativo: segurança semântica não
implica seletividade operacional, e a mesma união de predicados que elimina a
seletividade também elimina o poder de teste --- no plano global, nenhuma
confirmação da campanha dependia só da regra testada. O segundo é que a
garantia de um predicado derivado de regra depende de como o IDS avalia a
regra, e não só do texto do cabeçalho: a compilação que mantinha as portas
deixava escapar alertas no plano forte, e sua aparente seletividade vinha em
boa parte dessas portas. O terceiro é construtivo, mas limitado: a
proveniência indica quais regras importam para cada pacote, porém, com a
semântica do Snort~3, poucas regras restam discriminadas pelo cabeçalho, e a
redução sobre o ruleset fica em 18,3\%. Um seletor de candidatas que usasse
também o serviço identificado --- que o próprio Snort usa para rejeitar regras
de outros serviços --- é o passo natural seguinte e não foi avaliado aqui.
Todas as medidas de conjunto candidato são de viabilidade \emph{offline}: não
foi implementada execução do Snort orientada por conjunto de candidatos.

\subsection{Limitações}

Os resultados valem para um ruleset, um perfil de variáveis e um conjunto de
dados antigo em relação ao tráfego atual. Para as quantidades por pacote,
medimos também, offline, a quinta-feira inteira (9.322.025 quadros, 9.240.405
em $D$): o pré-filtro global continua em 100\%, o discriminador forte marca
95,335\%, e a redução média é de 99,42\% na partição elegível e 18,29\% sobre
o ruleset; as janelas não são atípicas do dia, mas os alertas do Snort não
foram calculados para o restante dele. A campanha amplia a cobertura de
regras para 958 e 134 confirmadas, mas com tráfego sintético, um ponto por
regra e sempre numa porta do próprio cabeçalho, de modo que só o teste
dedicado observa a avaliação service-based; as 1.875 regras sem testemunha
sintetizável não foram testadas. A semântica service-based foi estabelecida
pela leitura do código-fonte e da documentação do Snort~3.12.2.0 e confirmada
num teste de plano de dados só para um serviço e uma porta. A partição do discriminador forte não
é fechada sob flowbits: 24 arestas a cruzam, todas no flowbit
\texttt{smb.trans2}, e uma execução que separasse os dois lados teria de
manter esse componente junto. O alvo é o BMv2 em software: não houve medição
de vazão, latência ou recursos de hardware. Fora do ruleset avaliado, o
compilador só foi exercitado na compilação, sobre as regras do
Snort~2.9.15.1 distribuídas pelo Ubuntu; as dez regras só de cabeçalho desse
ruleset, únicas classificadas como \cls{Exact}, concordaram com o Snort nos
340 pares (quadro, regra) amostrados. Por fim, algumas condições de cabeçalho representáveis ainda são
omitidas, o que só alarga os predicados: faixas de porta, tipo e código ICMP
e \texttt{ip\_proto}. Representá-las elevaria a redução sobre o ruleset do
plano forte em no máximo 1,3, 3,1 e 0,07 ponto percentual, respectivamente.

\section{Conclusão e Trabalhos Futuros}

Sistemas de detecção de intrusão baseados em assinatura enfrentam custo
crescente de inspeção em software, e a pré-filtragem em switches programáveis
é a resposta natural. A literatura demonstrou a viabilidade dessa abordagem,
mas tratou a tradução parcial de regras como detalhe de implementação: quando
uma condição de payload ou de estado não cabe no plano de dados, o que foi
feito com ela não é explicitado. Isso é um problema mal resolvido porque o
erro resultante é silencioso --- um predicado mais restritivo que a regra de
origem descarta tráfego que deveria ser inspecionado, sem que nenhuma métrica
de cobertura de tradução o revele.

Este trabalho desenvolveu um compilador Snort-para-P4 que busca garantir que
todo pacote sobre o qual uma regra dispararia seja marcado como candidato, e
que registra explicitamente cada componente semântico omitido. Das 4.017
regras avaliadas, 3.996 resultaram em predicados necessários e 21 foram
declaradas não suportadas; o backend instalou 235 entradas em um alvo BMv2,
com fail-open fora do domínio que o parser interpreta. Os alertas observados
foram preservados, e a campanha fim a fim sobre 1.000 regras não encontrou
violação. O principal resultado é negativo e informativo: o pré-filtro
binário marcou 100\% dos pacotes como candidatos. A proveniência retida reduz
o conjunto de regras candidatas a cerca de quatro das 739 elegíveis, mas só
18,3\% sobre o ruleset; a compilação anterior, que mantinha portas que o
Snort~3 ignora em 2.931 regras service-based, prometia 81,8\% e deixava
escapar alertas. A contribuição para a literatura é a separação explícita
entre segurança semântica e seletividade operacional, acompanhada de um
manifesto de compilação que torna auditável o que cada tradução preserva, e a
constatação de que essa garantia, e a seletividade aparente, dependem de como
o IDS avalia a regra, e não apenas do seu texto.

Como trabalhos futuros, três direções se destacam. A primeira é um seletor de
regras candidatas que use o serviço identificado, além do cabeçalho, e a
execução do Snort orientada por esse conjunto, transformando a redução medida
\emph{offline} em ganho efetivo --- o que exige tratar as dependências de
estado entre regras, como as dos ACKs HTTP e dos flowbits da campanha. A
segunda é ampliar a expressividade do backend, incorporando faixas de porta
por decomposição ternária e campos ICMP à chave da tabela. A terceira é
avaliar o custo de execução em alvos de hardware e sobre conjuntos de regras,
perfis e tráfegos mais amplos, incluindo medição de vazão e latência, que não
foram objeto desta avaliação.

\bibliographystyle{IEEEtran}
\bibliography{references}

\end{document}
